\documentclass[conference]{IEEEtran}

\ifCLASSINFOpdf
\usepackage{amsmath}
\usepackage{amssymb}
\usepackage{booktabs}
\usepackage{tabularx}
\usepackage{graphicx}
\usepackage{float}
\usepackage{url}
\usepackage{enumitem}
\fi

\hyphenation{op-tical net-works semi-conduc-tor}

\newcommand{\missingfigurebox}[1]{%
\fbox{%
\begin{minipage}[c][0.22\textheight][c]{0.92\linewidth}
\centering\footnotesize Figure file not found:\\\texttt{\detokenize{#1}}
\end{minipage}}}

\newcommand{\reportimage}[2][\linewidth]{%
\IfFileExists{#2}{\includegraphics[width=#1]{#2}}{\missingfigurebox{#2}}}

\begin{document}

\title{Deep Saliency Map Prediction for Human Eye Fixation Using FCN-ResNet50}

\author{%
\begin{tabular}{c@{\hspace{0.9cm}}c@{\hspace{0.9cm}}c}
Nasrul Huda & Muhammad Bassam & Diksha Prasad\\
M.Sc. DSAI & M.Sc. IAS & M.Sc. DSAI\\
University of Hamburg & University of Hamburg & University of Hamburg
\end{tabular}
}

\maketitle

\begin{abstract}
This report presents a PyTorch-based eye fixation prediction system that maps natural images to dense $224\times224$ saliency maps. The system uses a fully convolutional network with an ImageNet-pretrained ResNet-50 backbone, followed by Gaussian smoothing and a log center-bias prior. Training first uses a frozen backbone and then end-to-end fine-tuning, improving the best validation loss from 0.189 to 0.173. The fine-tuned run also reaches a peak mean validation ROC-AUC of 0.922. Qualitative analysis of validation fixation maps shows that single-object scenes provide concentrated targets, while multi-person scenes contain distributed attention patterns that are more difficult to model.
\end{abstract}

\IEEEpeerreviewmaketitle

\section{Introduction}

\subsection{Objective}
The objective of this project was to develop a deep learning system using PyTorch that predicts human eye fixation maps from natural images, producing a dense $224\times224$ saliency map for each input image that indicates where viewers are likely to look.

\subsection{Approach Overview}
We implemented a fully convolutional network (FCN) based on a pretrained ResNet-50 backbone with a single-channel output head. Raw network logits are post-processed with a fixed Gaussian smoothing kernel and a log center-bias prior to better match the spatial smoothness and central tendency of human fixations. The model was first trained with a frozen backbone and subsequently fine-tuned end-to-end. Training used binary cross-entropy loss on pixel-wise fixation density maps, with ImageNet normalization applied to input images.

\section{System Architecture and Methodology}

\subsection{Foundational Work}
Our implementation follows the FCN paradigm for dense prediction introduced by Long et al.~\cite{long2015fcn}, adapted from semantic segmentation to fixation and saliency map prediction as outlined in the course exercises~\cite{course2026}. The encoder uses a ResNet-50 backbone pretrained on ImageNet~\cite{he2016resnet}, accessed via PyTorch's \texttt{fcn\_resnet50} implementation. Post-processing with a Gaussian kernel and center bias follows standard practice in computational saliency modeling to enforce spatial smoothness and account for the well-documented tendency of observers to fixate near image centers.

\subsection{Network Design}
\textbf{Backbone and decoder.} \texttt{FixationNet} uses torchvision's FCN-ResNet50 implementation with ImageNet-pretrained ResNet-50 backbone weights. The FCN head outputs a single channel of per-pixel logits at the input resolution ($224\times224$), treating each pixel as an independent Bernoulli target rather than a class label in a softmax sense.

\textbf{Post-processing modules.} Two non-trainable components refine the raw output:
\begin{enumerate}[leftmargin=*]
    \item \textbf{Gaussian smoothing:} A separable $25\times25$ Gaussian kernel ($\sigma=11.2$) is applied via \texttt{conv2d} with ``same'' padding, producing spatially smooth saliency maps consistent with aggregated eye-tracking data.
    \item \textbf{Center bias:} A dataset-provided center-bias density map, \texttt{center\_bias\_density.npy}, is stored in log-space and added to the smoothed logits, encoding the prior that fixations tend to fall near the image center.
\end{enumerate}

The forward pass is summarized as
\[
\begin{aligned}
&\text{FCN logits} \rightarrow \text{Gaussian smoothing}\\
&\rightarrow \text{log center bias} \rightarrow \text{output logits},
\end{aligned}
\]
with a sigmoid applied at inference.

\subsection{Modifications and Rationale}
We experimented with two training regimes, summarized in Table~\ref{tab:training_regimes}.

\begin{table}[H]
\centering
\small
\setlength{\tabcolsep}{3pt}
\renewcommand{\arraystretch}{1.1}
\begin{tabular}{@{}l l c c@{}}
\toprule
\textbf{Regime} & \textbf{Trainable layers} & \textbf{Best loss} & \textbf{Best/total} \\
\midrule
Frozen & FCN head & 0.189 & 16 / 88 \\
Fine-tuned & Full network & \textbf{0.173} & 24 / 38 \\
\bottomrule
\end{tabular}
\caption{Training regimes compared during development.}
\label{tab:training_regimes}
\end{table}

Initially, the ResNet-50 backbone was frozen to stabilize early training and reduce overfitting on the limited dataset. After observing a validation plateau around 0.19, we unfroze the full network and continued training end-to-end. Fine-tuning improved the best validation loss by approximately 8.3\% (0.189 to 0.173). The peak validation ROC-AUC also increased from 0.895 in the frozen run to \textbf{0.922} in the fine-tuned run, indicating that task-specific adaptation of low-level and mid-level features helps capture image-dependent saliency beyond what the generic center bias and frozen features can provide.

No additional architectural changes, such as attention modules or skip-connection modifications, were introduced. This kept the model aligned with the course template while focusing optimization effort on the training strategy.

\section{Training}

\subsection{Dataset and Preprocessing}
\textbf{Splits.} We used the course-provided fixation dataset with the official split files:
\begin{itemize}[leftmargin=*]
    \item \textbf{Training:} 3,006 image--fixation pairs.
    \item \textbf{Validation:} 1,128 image--fixation pairs.
    \item \textbf{Test:} 1,032 images, listed in \texttt{test\_images.txt}, used for final prediction submission.
\end{itemize}

\textbf{Input format.} Images are RGB at $224\times224$ pixels. Ground-truth fixation maps are single-channel grayscale density maps at $224\times224$, with pixel values in $[0,255]$ converted to $[0,1]$ via \texttt{ToTensor()}.

\textbf{Augmentation.} No data augmentation was applied in the final pipeline. Images were only converted to tensors and normalized. This choice keeps preprocessing simple and ensures validation metrics remain directly comparable across epochs; future work could add horizontal flips or color jitter to improve generalization.

\textbf{Normalization.} Input images were normalized with ImageNet statistics: mean = $[0.485, 0.456, 0.406]$ and standard deviation = $[0.229, 0.224, 0.225]$. Fixation maps received only \texttt{ToTensor()} scaling to $[0,1]$.

\subsection{Loss Function}
We trained with \textbf{binary cross-entropy with logits}, implemented as \texttt{F.binary\_cross\_entropy\_with\_logits}. This loss compares the model's per-pixel logits to the ground-truth fixation density after scaling to $[0,1]$. It is well suited for dense fixation prediction because each pixel is treated as an independent probability target, allowing the network to learn both high-saliency regions and background without requiring a normalized probability distribution over the image.

\subsection{Hyperparameters}
The final model was selected using the lowest validation loss observed during fine-tuning. The training configuration is shown in Table~\ref{tab:hyperparameters}.

\begin{table}[H]
\centering
\small
\renewcommand{\arraystretch}{1.1}
\begin{tabularx}{\linewidth}{@{}l X@{}}
\toprule
\textbf{Hyperparameter} & \textbf{Value} \\
\midrule
Optimizer & Adam \\
Learning rate & $1\times10^{-4}$ \\
Training batch size & 16 \\
Validation batch size & 64 \\
Dropout & None \\
Epochs trained (fine-tuning) & 38, best at epoch 24 \\
Total planned epochs & 100, with final selection by best validation loss \\
\bottomrule
\end{tabularx}
\caption{Final fine-tuning hyperparameters.}
\label{tab:hyperparameters}
\end{table}

Checkpoints were saved whenever validation loss improved, allowing training to resume from the best saved model state.

\section{Results}

\subsection{Quantitative Results}
The best validation loss achieved during fine-tuning was \textbf{0.1730} at epoch 24. Mean validation ROC-AUC was computed per image by applying a sigmoid to the model logits and binarizing ground-truth fixation pixels at 0.5. The ROC-AUC at the best-loss epoch was 0.918, and the peak ROC-AUC over the fine-tuning run was \textbf{0.922} at epoch 38.

Training and validation loss curves over 38 fine-tuning epochs are shown in Fig.~\ref{fig:loss_curves}. Training loss decreased steadily from 0.319 to 0.125, while validation loss dropped sharply in the first few epochs and then plateaued around 0.173--0.177, suggesting mild overfitting in later epochs but stable generalization after the best checkpoint.

\begin{figure}[H]
\centering
\reportimage{report_figures/loss_curves.png}
\caption{Training and validation binary cross-entropy loss over 38 fine-tuning epochs.}
\label{fig:loss_curves}
\end{figure}

Fig.~\ref{fig:metric_comparison} compares the frozen-backbone and fine-tuned training histories. The frozen model improves quickly but remains near a validation-loss floor of 0.19, whereas the fine-tuned model consistently operates at lower validation loss and higher ROC-AUC.

\begin{figure}[H]
\centering
\reportimage{report_figures/metric_comparison.png}
\caption{Validation loss and ROC-AUC for the frozen-backbone and fine-tuned training regimes.}
\label{fig:metric_comparison}
\end{figure}

\begin{table}[H]
\centering
\small
\setlength{\tabcolsep}{4pt}
\renewcommand{\arraystretch}{1.1}
\begin{tabular}{@{}l c c@{}}
\toprule
\textbf{Metric} & \textbf{Frozen backbone} & \textbf{Fine-tuned} \\
\midrule
Val loss & 0.189 & \textbf{0.173} \\
Peak val ROC-AUC & 0.895 & \textbf{0.922} \\
Best-loss epoch & 16 & 24 \\
\bottomrule
\end{tabular}
\caption{Validation performance for the frozen and fine-tuned models.}
\label{tab:validation_results}
\end{table}

\subsection{Qualitative Analysis}
The available validation images and fixation maps illustrate two important target structures for this task: compact, single-object fixation maps and distributed, multi-object fixation maps.

\subsubsection{Concentrated case: \texttt{image-4026}}
Fig.~\ref{fig:case_4026} shows a validation image with a single person in an open landscape. The ground-truth fixation density is compact and centered on the person, making this type of example well matched to the model's Gaussian smoothing and center-bias assumptions.

\begin{figure}[H]
\centering
\reportimage{report_figures/case_4026.png}
\caption{Validation example with a compact fixation target. The fixation map is concentrated around the single visible person.}
\label{fig:case_4026}
\end{figure}

\subsubsection{Distributed case: \texttt{image-3829}}
Fig.~\ref{fig:case_3829} shows a multi-person golf-course scene. The ground-truth fixation map is spread across people, gestures, the golf bag, and background regions. Such examples are more challenging because the target is multi-modal and depends on semantic relationships between objects rather than only local contrast or centrality.

\begin{figure}[H]
\centering
\reportimage{report_figures/case_3829.png}
\caption{Validation example with distributed fixation targets. The target map contains several separated high-attention regions across people and contextual objects.}
\label{fig:case_3829}
\end{figure}

These examples explain the main qualitative limitation observed during development: the model is better suited to scenes with a single dominant salient object than to scenes where human attention is distributed across several semantic regions.

\section{Conclusion}
We built an eye fixation prediction system using an FCN-ResNet50 architecture with Gaussian smoothing and center-bias post-processing, trained on 3,006 images and validated on 1,128. Starting from a frozen ImageNet backbone with best validation loss 0.189, end-to-end fine-tuning yielded a best validation loss of \textbf{0.173} and a peak validation ROC-AUC of \textbf{0.922}. Qualitative inspection of validation targets confirms that scenes with clear, isolated salient objects are structurally easier than complex multi-person images with distributed fixation regions.

Key lessons include the value of fine-tuning over frozen features for this task and the importance of inspecting failure modes. Border artifacts and missed secondary fixations point to limitations of pixel-wise BCE without explicit handling of sparsity or multi-modal attention.

Future work should explore KL-divergence or correlation-based losses better suited to sparse, probabilistic fixation maps. Further improvements could also add data augmentation or an attention mechanism to improve performance on multi-object scenes.

\begin{thebibliography}{1}
\bibitem{long2015fcn}
J. Long, E. Shelhamer, and T. Darrell, ``Fully Convolutional Networks for Semantic Segmentation,'' \textit{Proceedings of the IEEE Conference on Computer Vision and Pattern Recognition (CVPR)}, 2015.

\bibitem{he2016resnet}
K. He, X. Zhang, S. Ren, and J. Sun, ``Deep Residual Learning for Image Recognition,'' \textit{Proceedings of the IEEE Conference on Computer Vision and Pattern Recognition (CVPR)}, 2016.

\bibitem{course2026}
Course materials, Computer Vision 2, Exercise 5 --- FCN for Eye Fixation Prediction, University of Hamburg, 2026.
\end{thebibliography}

\end{document}
